\section{Pure reasoning benchmarks：尽量剥离知识}

前面的任务大量依赖语言知识和世界知识。ARC-AGI 试图问一个更纯的问题：能否把 reasoning 从 memorized facts 中分离出来？它用网格变换任务让每题都相对独特，减少直接记忆帮助。

\begin{figure}[H]
\centering
\includegraphics[width=0.9\textwidth]{arc-task-grids.jpg}
\caption{ARC-AGI-1 grid tasks：每个任务独特，人类可解，记忆帮助有限。}
\figsource{ARC Prize 官方图，本地化后供 CS336 Lecture 12 使用。}
\end{figure}

\begin{knowledgebox}{读图：ARC-AGI 想测什么}
ARC-AGI 试图把推理从语言知识中剥离出来。输入输出是小网格转换，要求从少量 examples 推断规则。它的目标不是百科知识，而是抽象模式发现。读这些网格时，真正要看的是“从几个例子归纳规则，再应用到 test grid”的能力。
\end{knowledgebox}

\begin{figure}[H]
\centering
\includegraphics[width=0.9\textwidth]{arc-agi-2-unsolved.png}
\caption{ARC-AGI-2：更强调 multi-step reasoning。}
\figsource{ARC Prize 官方图，本地化后供 CS336 Lecture 12 使用。}
\end{figure}

\begin{figure}[H]
\centering
\includegraphics[width=0.86\textwidth]{arc-agi-results.png}
\caption{ARC-AGI results：pretrained LMs 提升有限，reasoning models 开始推动成绩。}
\figsource{CS336 Lecture 12 官方源引用图。}
\end{figure}

\begin{knowledgebox}{老师强调：reasoning benchmark 也有边界}
纯预训练语言模型不一定擅长搜索和多步假设检验。o1/o3 类 reasoning models 开始在 ARC 上提升，说明训练范式会改变曲线。但 ARC 仍是人类设计的抽象任务，它测的是某种受控 reasoning，不是所有现实推理。
\end{knowledgebox}

\begin{figure}[H]
\centering
\includegraphics[width=0.45\textwidth]{arc-agi-3.png}
\caption{ARC-AGI-3：转向 interactive environments。}
\figsource{CS336 Lecture 12 官方源引用图。}
\end{figure}

\begin{figure}[H]
\centering
\includegraphics[width=0.72\textwidth]{arc-agi-3-results.png}
\caption{ARC-AGI-3 results：交互式环境进一步测试模型适应和探索。}
\figsource{CS336 Lecture 12 官方源引用图。}
\end{figure}

\begin{warningbox}{纯 reasoning 也不是纯粹无争议}
ARC 试图减少知识记忆，但任务设计仍包含人类先验、视觉模式偏好和搜索空间选择。它测的是一种人类可理解的抽象推理，不代表所有 intelligence。
\end{warningbox}

